\documentclass[10pt,a4paper]{amsart}
\usepackage[margin=19mm]{geometry}
\usepackage{amsmath,amssymb,mathtools}
\usepackage[scr=rsfs,cal=euler]{mathalfa}
\usepackage{xfrac}
\usepackage[unicode,hidelinks,bookmarks=false]{hyperref}
\newcommand{\ie}{i.e.\ }
\newcommand{\cf}{cf.\ }
\newcommand{\resp}{respectively\ }
\newcommand{\Subsection}{\subsection}
\providecommand{\nobreakdash}{\nobreak\hbox{-}}
\setlength{\emergencystretch}{2em}
\multlinegap0pt
\usepackage{amssymb}
\usepackage{mathtools}
\usepackage{xcolor}
\usepackage[shortlabels]{enumitem}
\newtheorem{lemma}{Lemma}
\newtheorem{theorem}{Theorem}
\usepackage{cleveref}
\crefname{lemma}{Lemma}{Lemmas}
\crefname{theorem}{Theorem}{Theorems}
\newcommand{\R}{\mathbb{R}}
\newcommand{\abs}[1]{\lvert #1\rvert}
\newcommand{\froinner}[2]{\langle #1,#2\rangle_{\!F}}
\newcommand{\inner}[2]{\langle #1,#2\rangle}
\newcommand{\norm}[1]{\lVert #1\rVert}
\DeclareMathOperator{\tr}{tr}
\title{Upper bound on the $k$-th eigenvalue of a graph}
\author{Varun Sivashankar}
\date{Source: arXiv:2603.28738v1 (CC BY 4.0)}
\begin{document}
\maketitle

\noindent\emph{Source and licence.} Upper bound on the $k$-th eigenvalue of a graph. Version arXiv:2603.28738v1 (CC BY 4.0), distributed by arXiv under the Creative Commons Attribution 4.0 International licence, http://creativecommons.org/licenses/by/4.0/, as recorded in the arXiv licence metadata for this exact version and shown on the abstract page https://arxiv.org/abs/2603.28738v1. Full licence notice: \texttt{licenses/arxiv-2603.28738-CC-BY-4.0.txt}.\par\smallskip

\noindent\textit{Editorial scope.} Counted unit: the article's theorem thm:l1-proj (source0.tex lines 231--240). The article's complete original proof of it is delivered with it (source0.tex lines 241--421, one \texttt{proof} environment of 181 lines). No mathematics is written, repaired or re-derived: the statement and the proof are the article's own lines, in the article's own notation, and no retained line is substituted or re-flowed.  Retained uncounted as the source-local context the proof needs: lines 174--185; lines 203--209; lines 553--553.  Nothing was omitted from any retained span, so the package carries no omission record.  No retained line contains a citation, so the wrapper carries no bibliography and no bibliography entry.  No retained line contains a figure environment, an image-inclusion command or drawing code, so no image is delivered and the package declares no asset dependency.  Editorial formatting only, in addition: an AMS \texttt{amsart} preamble that restates the article's own declarations and macros used by the retained lines, namely the environments \texttt{lemma}, \texttt{theorem} on separate counters, together with the standard packages those lines need.  The retained environments print in this document as Lemma 1, Theorem 1 in order of appearance, whereas the article prints the same items with the numbering of its own full document.  The complete native source of the article as distributed in the arXiv e-print for this exact version is bundled unchanged as \texttt{sources/197466/source0.tex}.
\begin{lemma}\label{lem:psd2}
For $r \geq 2$, define
\[
\phi(u):=uu^\top-\frac1r I_r,
\qquad u\in S^{r-1}.
\]
Then for all $u,v\in S^{r-1}$,
\[
\froinner{\phi(u)}{\phi(v)}=f_2^r(\inner{u}{v}).
\]
Further, $f_2^r$ is positive semidefinite on $S^{r-1}$.
\end{lemma}

\begin{lemma}\label{lem:psd4}
For $r \geq 2$,
\[
f_4^r(t)=t^4-\frac{3}{r(r+2)}
\]
is positive semidefinite on \(S^{r-1}\).
\end{lemma}

\section{Concluding Remarks}\label{sec5}

\begin{theorem}\label{thm:l1-proj}
Let \(Q\in \R^{n\times n}\) be a rank-\(r\) orthogonal projection, where \(r\ge 2\). 
% Let
% \[\beta_r:=\frac{r+\sqrt{r+2}}{1+\sqrt{r+2}}\]
% Then
\[
\|Q\|_1=\sum_{i,j=1}^n \abs{q_{ij}}
\le \beta_r n \qquad \text{ where } \quad \beta_r:=\frac{r+\sqrt{r+2}}{1+\sqrt{r+2}}
\]
\end{theorem}

\begin{proof}
Since $Q$ is a rank $r$ orthogonal projection, we may write $Q = BB^\top$ such that $B \in \R^{n\times r}$ where the columns of $B$ form an orthogonal set in $\R^n$. Further, let $x_1,\ldots,x_n \in \R^r$ be the rows of $B$. Write $x_i=c_i u_i$, where $c_i=\norm{x_i}_2\ge 0$ and $u_i\in S^{r-1}$ whenever $c_i\ne 0$; if $c_i=0$, choose $u_i$ arbitrarily in $S^{r-1}$. Since the columns of $B$ are orthogonal, $B^\top B = I_r$. Writing out $B^\top B$ in terms of the rows of $B$, we obtain
\begin{equation}\label{eq:projection-model}
\sum_{i=1}^n c_i^2 u_i u_i^\top=I_r 
\qquad\text{and}\qquad
q_{ij}=c_i c_j\inner{u_i}{u_j}
\end{equation}


% In particular,
% \begin{equation}\label{eq:sum-ci2}
% \sum_{i=1}^n c_i^2=\tr(Q)=r.
% \end{equation}


Set
\[
C:=\sum_{i=1}^n c_i,
\qquad
x:=\sum_{i=1}^n c_i\phi(u_i),
\qquad
X:=\norm{x}_F.
\]

To prove \Cref{thm:l1-proj}, we establish \Cref{lem:CS-ineq}, \Cref{lem:abstract-majorant} and \Cref{lem:scalar-majorant}.

\begin{lemma}\label{lem:CS-ineq}
With the notation above,
\[
C^2+rX^2\le rn.
\]
\end{lemma}

\begin{proof}
\begin{align*}
C^2 + rX^2
&= \sum_{i=1}^n \sum_{j=1}^n c_i c_j + r \sum_{i=1}^n \sum_{j=1}^n c_i c_j \froinner{\phi(u_i)}{\phi(u_j)}\\
&= \sum_{i=1}^n \sum_{j=1}^n c_i c_j + r \sum_{i=1}^n \sum_{j=1}^n c_i c_j \left(\inner{u_i}{u_j}^2 - \frac{1}{r}\right) \text{\qquad by \Cref{lem:psd2}}\\
&= r \sum_{i=1}^n \sum_{j=1}^n c_i c_j \inner{u_i}{u_j}^2\\
&\leq \frac{r}{2}\sum_{i=1}^n \sum_{j=1}^n (c_i^2 + c_j^2) \inner{u_i}{u_j}^2 \text{\qquad by Cauchy-Schwarz}\\
&= r \sum_{i=1}^n \sum_{j=1}^n c_j^2 \inner{u_i}{u_j}^2 \text{\qquad by symmetry}
% &= r \sum_{i=1}^n \sum_{j=1}^n \tr(u_i u_i^\top c_j^2 u_j u_j^\top)\\
% &= r \sum_{i=1}^n  \tr\left(u_i u_i^\top \sum_{j=1}^n c_j^2 u_j u_j^\top\right)\\
% &= r \sum_{i=1}^n  \tr\left(u_i u_i^\top\right) \text{\qquad by \Cref{eq:projection-model}}\\
% &= r n
\end{align*}

\begin{align*}
r \sum_{i=1}^n \sum_{j=1}^n c_j^2 \inner{u_i}{u_j}^2
&= r \sum_{i=1}^n \sum_{j=1}^n \tr(u_i u_i^\top c_j^2 u_j u_j^\top)\\
&= r \sum_{i=1}^n  \tr\left(u_i u_i^\top \sum_{j=1}^n c_j^2 u_j u_j^\top\right)\\
&= r \sum_{i=1}^n  \tr\left(u_i u_i^\top\right) \text{\qquad by \Cref{eq:projection-model}}\\
&= r n
\end{align*}
\end{proof}



We will now bound \(\|Q\|_1\) by finding a polynomial $h(t)$ that satisfies $|t| \leq h(t)$ for all $t \in [-1,1]$ and applying it entry wise to $|q_{ij}| = c_i c_j |\inner{u_i}{u_j}|$ with $t = |\inner{u_i}{u_j}|$. We will consider an expression of the form
\[
a+b\,f_2^r(t)-\gamma\,f_4^r(t).
\]
This is because after substituting \(t=\inner{u_i}{u_j}\) and summing over \(i,j\), the constant term produces $C^2$, the \(f_2^r\)-term produces \(X^2\), while the \(f_4^r\)-term is nonnegative since $f_4^r$ is positive semidefinite by \Cref{lem:psd4}. The following lemma formalizes this reduction.

\begin{lemma}\label{lem:abstract-majorant}
Let \(a,b,\gamma\in\R\) satisfy
\[
\abs{t}\le a+b\,f_2^r(t)-\gamma\,f_4^r(t)
\qquad\text{for all } t\in[-1,1],
\]
with \(\gamma\ge 0\) and \(b\le ra\). Then every rank-\(r\) orthogonal projection \(Q\in\R^{n\times n}\) satisfies
\[
\|Q\|_1\le arn.
\]
\end{lemma}

\begin{proof}
Recall by \eqref{eq:projection-model} that we may write $Q = BB^T$ where the rows of $B$ are given by $c_i u_i$ for $i \in [n]$ and $\|Q\|_1 = \sum_{i,j} |q_{ij}| = \sum_{i,j} c_i c_j |\inner{u_i}{u_j}|$. Applying the assumed scalar inequality with \(t=\inner{u_i}{u_j}\), using the fact that $c_i \geq 0$ and summing over all \(i,j\), we obtain
\[
\|Q\|_1 = \sum_{i,j} |q_j|
\le a\sum_{i,j} c_i c_j
+b\sum_{i,j} c_i c_j f_2^r(\inner{u_i}{u_j})
-\gamma\sum_{i,j} c_i c_j f_4^r(\inner{u_i}{u_j}).
\]
By definition,
\[
\sum_{i,j} c_i c_j=C^2,
\qquad
\sum_{i,j} c_i c_j f_2^r(\inner{u_i}{u_j}) = \sum_{i,j} c_i c_j \froinner{\phi(u_i)}{\phi(u_j)} = X^2
\]
Also, the positive semidefiniteness of \(\bigl(f_4^r(\inner{u_i}{u_j})\bigr)_{i,j=1}^n\) implies that
\[
\sum_{i,j} c_i c_j f_4^r(\inner{u_i}{u_j})\ge 0.
\]
Therefore
\[
\|Q\|_1\le aC^2+bX^2.
\]
Since \(b\le ra\), we have
\[
aC^2+bX^2\le a\bigl(C^2+rX^2\bigr).
\]
Now apply \Cref{lem:CS-ineq} to get
\[
\|Q\|_1\le arn.
\]
\end{proof}

We now exhibit a concrete choice of coefficients for which the hypotheses of \Cref{lem:abstract-majorant} hold. Clearly, we would like $ar$ to be as small as possible. The choice of coefficients in \Cref{lem:scalar-majorant} below may seem mysterious but it is motivated by a connection to equiangular lines. A more detailed discussion is included in \Cref{sec5}. 

% We also checked \Cref{lem:scalar-majorant} using Lean.  

\begin{lemma}\label{lem:scalar-majorant}
Let \(r\ge 2\), and set
\[
s:=\sqrt{r+2}.
\]
Define
\[
a_r:=\frac{s^2+s-2}{r(s+1)},\qquad
b_r:=\frac{s(s^2+2s+3)}{2(s+1)^2},\qquad
\gamma_r:=\frac{s^3}{2(s+1)^2}.
\]
Then for every \(t\in[-1,1]\),
\[
|t|\le a_r+b_r f_2^r(t)-\gamma_r f_4^r(t).
\]
Moreover,
\[
b_r\le ra_r.
\]
\end{lemma}

\begin{proof}
Recall that
\[
f_2^r(t)=t^2-\frac1r,\qquad
f_4^r(t)=t^4-\frac{3}{r(r+2)}.
\]
Since the left-hand side is even in \(t\), it suffices to consider \(t\in[0,1]\).
A direct expansion gives
\[
a_r+b_r f_2^r(t)-\gamma_r f_4^r(t)-t
=
\frac{(1-t)(st-1)^2(st+s+2)}{2(s+1)^2}.
\]
The right-hand side is nonnegative for \(t\in[0,1]\), since
\[
1-t\ge 0,\qquad (st-1)^2\ge 0,\qquad st+s+2>0.
\]
Hence
\[
|t|\le a_r+b_r f_2^r(t)-\gamma_r f_4^r(t)
\qquad\text{for all } t\in[-1,1].
\]

Finally,
\[
ra_r-b_r
=
\frac{s^3+2s^2-5s-4}{2(s+1)^2}
=
\frac{(s-2)(s+1)(s+3)+2}{2(s+1)^2}\ge 0,
\]
and therefore \(b_r\le ra_r\).
\end{proof}



\noindent\textit{Completing the proof of \cref{thm:l1-proj}}.

Apply \Cref{lem:abstract-majorant} with
\[
a=a_r,\qquad b=b_r,\qquad \gamma=\gamma_r.
\]
By \Cref{lem:scalar-majorant}, the hypotheses of \Cref{lem:abstract-majorant} are satisfied. Therefore
\[
\|Q\|_1\le a_r r n = \frac{s^2 + s -2}{s+1} n = \frac{r+\sqrt{r+2}}{1+\sqrt{r+2}} n = \beta_r n.
\]

\end{proof}

\end{document}
